\section{ARC-01--05: repairing the archive argument mathematically}
The relevant theorem on cote 1 is the Hilbert trace/determinant theorem on scans 79--80, with the coefficient domination on scan 79 and the genus-zero argument on scans 111--112. Some of the trace theorem text is struck out. We do not turn that editorial status into a modern publication-priority assertion. The purpose here is to check the mathematical mechanism and specify repairs to the supplied text. Original page images would be needed to call a repair a certified transcription correction.

\subsection{ARC-01: the complete corrected genus-zero lemma}
\begin{lemma}
Let $f(z)=\sum a_nz^n$ be a nonzero entire function and $h(z)=\sum b_nz^n$ an entire function of genus zero with $b_n\ge0$ and $|a_n|\le b_n$. If the nonzero zeros $z_j$ of $f$, counted with multiplicity, satisfy $\sum_j|z_j|^{-1}<\infty$, then $f$ has genus zero.
\end{lemma}
\begin{proof}
A genus-zero product $Cz^m\prod_j(1-z/u_j)$ with $\sum|u_j|^{-1}<\infty$ has minimal exponential type. Indeed, for any $\varepsilon>0$, choose a tail with reciprocal sum less than $\varepsilon/2$. Its modulus is at most $e^{\varepsilon|z|/2}$, while the remaining polynomial is at most $C_\varepsilon e^{\varepsilon|z|/2}$. Thus
\[
 |f(z)|\le h(|z|)\le C_\varepsilon e^{\varepsilon|z|}.
\]
Hadamard factorization \cite[Theorem 3.16, p.~92]{mcmullen} at order at most one and the assumed reciprocal summability give
$f(z)=Ce^{bz}g(z)$ with $C\ne0$ and $g(z)=z^m\prod_j(1-z/z_j)$ of genus zero. If factorization is initially written with genus-one factors, the absolutely convergent sum $\sum1/z_j$ is absorbed into $b$. This standard factorization step is the background theorem invoked on scan 111, line 13.

If $b\ne0$, replace $z$ by $(\bar b/|b|)z$ so that the new exponent is $|b|z$. Choose $\pi/3<\theta<\pi/2$. On $|\arg z|\le\theta$ one has $\Re z\ge\cos\theta\,|z|$. Applying the preceding bound for $f$ with $\varepsilon=|b|\cos\theta/2$ yields $|g(z)|\le A e^{-a|z|}$ on that sector, where $a>0$. The genus-zero bound for $g$ gives $|g(z)|\le B e^{a|z|/3}$ everywhere.

Let $\omega=e^{2\pi i/3}$ and $G_0(z)=g(z)g(\omega z)g(\omega^2z)$. For every direction one of the three rotated arguments has argument in $[-\pi/3,\pi/3]$. All have equal modulus. Therefore $|G_0(z)|\le AB^2e^{-a|z|/3}$ everywhere. Liouville's theorem \cite[Theorem 1.18, p.~19]{mcmullen} makes this entire function constant, and the decay makes the constant zero. A product of three nonzero entire functions is not identically zero. This contradiction forces $b=0$ and proves the lemma.
\end{proof}
If the initial genus-zero majorant has arbitrary complex coefficients, coefficientwise absolute values can be majorized by
$|C|(1+z)^m\prod_j(1+z/|u_j|)$; this justifies the positive-coefficient reduction. The required majorization is radial, $h(|z|)$, not $|h(z)|$. For instance $f=1-z$, $h=1+z$, $z=-1$ disproves the latter inequality. The result applies to arbitrary trace-class Hilbert operators via singular-value determinant domination, not through a positivity hypothesis on the original operator.

\subsection{ARC-02: the order in scan 110}
For $\alpha>1$, $P_\alpha(z)=\prod_{n\ge1}(1-z/n^\alpha)$ has order $1/\alpha$, not $1$.
\begin{proof}
Reciprocal summability gives genus zero. Let $N=\lfloor r^{1/\alpha}\rfloor$. At $z=-r$ the first $N$ factors are at least two, hence $\log M(r)\ge N\log2$. Conversely
\[
 \log M(r)\le \sum_{n\le N}\log(1+r/n^\alpha)+r\sum_{n>N}n^{-\alpha}.
\]
The first sum is at most $N\log(2r)-\alpha\log(N!)=O_\alpha(N)$, using $\log(N!)\ge N\log N-N$ and $r<(N+1)^\alpha$. The integral estimate for the tail also gives $O_\alpha(N)$. Thus $\log M(r)$ is bounded above and below by positive multiples of $r^{1/\alpha}$, proving the order.
\end{proof}
The source says ``sauf erreur,'' so an authorial tentative error is a real possibility; mathematical certainty does not identify the mistaken glyph.

\subsection{ARC-03--04: scan 113 requires a liminf and a prefactor}
For a positive nonincreasing sequence $a_n\to0$ define
\[
 r=\inf\{s>0:a_n=O(n^{-1/s})\},\qquad
 \sigma=\inf\{s>0:\sum_na_n^s<\infty\}.
\]
Then
\begin{equation}\label{eq:liminf}
 \frac1r=\liminf_{n\to\infty}\frac{\log(1/a_n)}{\log n},\qquad r=\sigma,
\end{equation}
with extended reciprocal conventions.
\begin{proof}
If $a_n\le Cn^{-1/s}$, the liminf in \eqref{eq:liminf} is at least $1/s$. Conversely if $1/s$ is strictly less than that liminf, the required bound holds for all sufficiently large $n$ and extends to the finite prefix by a constant. Taking thresholds proves the first assertion. If $a_n=O(n^{-1/s})$, then $\sum a_n^t$ converges for every $t>s$, so $\sigma\le r$. If $\sum a_n^s<\infty$, monotonicity gives
\[
 \lfloor n/2\rfloor a_n^s\le\sum_{j=\lceil n/2\rceil}^n a_j^s\longrightarrow0.
\]
Thus $a_n=o(n^{-1/s})$, proving $r\le\sigma$ and the little-$o$ assertion of the manuscript.
\end{proof}
A counterexample to the printed limsup is explicit. Let $N_k=2^{2^k}$ for $k\ge0$, $a_1=1$, and $a_n=N_k^{-1}$ for $N_k\le n<N_{k+1}=N_k^2$. The sequence decreases to zero. The logarithmic ratio is one at $N_k$ and tends to $1/2$ at the right ends of the blocks, so its limsup is one and liminf is $1/2$. Also $a_n\le n^{-1/2}$, while $a_nn^{1/s}$ is unbounded along block ends for $s<2$. Thus $r=2$, refuting the asserted limsup formula.

For the manuscript's alternative proof that $\sum_n a_n<\infty$ implies $na_n\to0$, suppose $a_{n_i}\ge1/n_i$ after a harmless rescaling and $\varepsilon_i=(n_i-n_{i-1})/n_i$, then for $i\ge2$ summability gives $\sum\varepsilon_i<\infty$, and the correct recurrence is
\[
 n_i=\frac{n_1}{\prod_{j=2}^i(1-\varepsilon_j)}
 \le\frac{n_1}{\prod_{j=2}^\infty(1-\varepsilon_j)}<\infty.
\]
The infinite product is positive because eventually $-\log(1-\varepsilon_j)\le2\varepsilon_j$. This contradicts strict increase of the integers. The prefactor is necessary unless an initial value one is explicitly fixed; an initial zero would instead make the first factor vanish. To prove little-$o$, the subsequence threshold is an arbitrary positive constant, reduced to one by rescaling.

\subsection{ARC-05: norm and domain readings that the formulas force}
For a compact Hilbert operator with singular values $s_j$, the nonzero singular values of $\Lambda^n u$ are the products $s_{i_1}\cdots s_{i_n}$ over strictly increasing indices. This follows first for finite-rank singular-value decompositions and then by finite-rank approximation. For $u$ trace class in the first identity and $u$ Hilbert--Schmidt in the second, this gives the finite identities
\[
 \|\Lambda^n u\|_1=\sum_{i_1<\cdots<i_n}s_{i_1}\cdots s_{i_n}
 =\alpha_n((u^*u)^{1/2}),\qquad
 \|\Lambda^n u\|_2^2=\alpha_n(u^*u).
\]
This proves the readings $\sqrt{u^*u}$ in scan 78, (16), and subscript $2$ in (17). At $n=1$, $u=I_2$, the latter square-root expression is $\sqrt2$, whereas $\|I_2\|_1=2$, so reading subscript $1$ is concretely false. The mixed bound (15) uses nuclear norms; at $n=1$, $u=I_N$, operator norm would give the false inequality $N\le1$. This is a norm-convention clarification, not proof that an omitted subscript was present in the manuscript.

In scan 81, $f$ is defined on the open convex set $U$. Its supporting affine functions satisfy $L_a(x)\le f(x)$ for $x\in U$ and $L_a(a)=f(a)$. For $x\in U$, choosing $a=x$ proves $f(x)=\sup_aL_a(x)$ with attainment. Without an extension the expression $f(x)$ for $x\notin U$ is undefined. Even the natural extension $F=\sup_aL_a$ need not attain its supremum outside $U$: on $U=(0,\infty)$, take $f(x)=1/x$, whose tangents have $L_a(0)=2/a$. Then $F(0)=+\infty$ and no tangent attains it. The required domain in item 3 is therefore $x\in U$, unless an extension and extra attainment hypotheses are supplied.

The expression $e^{z^\omega}$ on scan 110, line 1, is an entire example with the stated order and genus when $\omega$ is a nonnegative integer. For an integer $m>0$, $M(r)=e^{r^m}$ gives order $m$, and the exponential polynomial has genus $m$; $m=0$ gives a nonzero constant of order and genus zero. A nonintegral power has nontrivial monodromy near zero, which exponentiation cannot cancel identically on a punctured neighborhood; a negative integral power gives an essential singularity at zero. Thus the mathematical example needs that qualification, though it may have been implicit in its context.

\subsection{What the correction table does and does not assert}
The table records the product sign on 111/9, rotation on 111/16, sector angle on 112/3, radial majorant on 112/7 and its margin, third root on 112/17, order on 110/2--3, liminf on 113/8 and product indices on 113/22. It also records the unambiguous Hilbert--Schmidt norm on 78/(17), the square root $\sqrt{u^*u}$ on 78/(16), and the domain qualification on scan 81. Contextual issues such as the multiplicative constant and the absorbed $\beta$ factor in 112/6 are explicitly marked as possible implicit normalizations. We do not force a choice between $\sqrt{uu^*}$ and $\sqrt{u^*u}$ in scan 79: their relevant nonzero singular-value data agree, so this mathematics does not resolve that reading. Nor do we call the source's word ``constant'' wrong merely because exponential decay yields the stronger contradiction of a zero constant.

\section{Where these findings live in Grothendieck's frameworks}
\begin{longtable}{>{\raggedright\arraybackslash}p{0.23\textwidth}>{\raggedright\arraybackslash}p{0.40\textwidth}>{\raggedright\arraybackslash}p{0.27\textwidth}}
\toprule
Findings & Exact textual home / comparison & Scope of the relationship\\\midrule\endhead
B-01--03, DET-01--05, DET-08 & Cote 1, scans 78--80, (15)--(23); scans 111--112. SGA 6, V, 2.1, 2.2(b), 2.2.1, (7.1.1). & Analytic trace-class limits and finite exterior-power identities; no arithmetic realization supplied.\\
DET-06--07, DET-10, STEP7-01--02 & Cote 1/79, comparison with $|u|$; SGA 6, V, 2.2(b). & Finite period/Gram/Capelli determinants; a signed determinant is not a positive metric.\\
DET-12--17, DET-19 & Cote 1/78--80, trace ideals and Fredholm determinant; SGA 6, V, 2.1, 2.2(b), 2.2.1, (7.1.1). & Positive determinant lines, frames, Bessel convolutions and arithmetic-kernel reconstruction. The latter uses an RH-equivalent sign criterion. The natural boost flow is ruled out; historical finite matching remains separate.\\
DET-09, DET-11, DET-18, STEP6-01--03 & Cote 1/79--80, trace and resolvent; SGA 6, V, (7.1.1), for integral multiplicities; SGA 5, III B, 1.3.1 as the curve analogue. & Exact scalar prime-sum and Weil-form traces; rational family weights give a scaled ordinary trace after integer multiplication. Historical matrix matching remains separate.\\
TOY-01--02 & SGA 5, III, 4.7--4.8, for a realized sheaf trace; cote 1/79 for the distinct determinant trace. & Finite-field schemes and character sheaves when the weight has that realization; arbitrary weights need no sheaf.\\
MAS-01--12, MAS-BOUND & Cote 1/79--80 and 111--112; SGA 5, III, 4.7 only as geometric analogy. & The spectral quotient claim gives ordinary zeta and paired Hecke determinants. Their inverse-square operators are not positive in the actual quotient metrics. Draft dependence remains; no \'etale cohomology or crystal is identified.\\
PRG-WEIL, PRG-LI, PRG-NB, PRG-DB & Cote 1/79--80 for applicable operator identities; SGA 5, III B, 1.3.1 supplies the curve comparison. & The criteria themselves belong to their cited analytic works; no invented EGA/FGA locator.\\
PRG-CONNES, PRG-DEN & SGA 5, III, 4.7--4.8 and III B, 1.3.1; cote 1/79--80 separates ordinary from regularized determinants. & Proposed arithmetic analogues require their own constructions and adjoint/duality structure.\\
PRG-FF & SGA 5, III B, 1.3.1; Weil I, (1.5.1), (1.5.3)--(1.5.4), (1.6). & Actual $\ell$-adic cohomology, weights, finite determinants and the correct Frobenius convention.\\
ARC-01--05 & Cote 1, scans 78, 81, 110--113, individual lines in the transcription table. & Analytic repairs established; manuscript readings await the images.\\\bottomrule
\end{longtable}
An exact EGA or FGA citation cannot honestly locate a number-field Hilbert positivity construction that those texts do not make. Likewise this prompt produces no new theorem about crystals or divided powers. The operative common language here is exterior powers, traces, determinants, duality and weights. In schemes over $\Z$, local purity of an actual $\ell$-adic sheaf and global RH for its completed $L$-function are distinct assertions; this audit does not extend the stipulated degree-one results to general Hasse--Weil $L$-functions.
